\chapter{C++23 --- The Completion Release}
\label{ch:cpp23-completion-release}

C++20 was a revolution: concepts, ranges, coroutines, and modules arrived together and changed the shape of the language. But a revolution leaves loose ends. C++23 is the release that ties them off. It is not defined by one headline feature; it is defined by \emph{closure} --- finishing the standard library's coroutine story with \lstinline|std::generator|, completing the ranges library that shipped half-built, giving error handling a first-class vocabulary with \lstinline|std::expected|, and making the day-one ergonomics (\lstinline|std::print|, \lstinline|import std;|, deducing \lstinline|this|) finally usable. This chapter frames the volume: what kind of release C++23 is, why ``ergonomics and completion'' is the right lens, and how the chapters ahead are organized.

\section{The Shape of a Completion Release}

The C++ standards committee has settled into a three-year cadence, and within that cadence the releases alternate in character. \textbf{C++11 and C++20 were \emph{major} releases} --- each redefined how idiomatic C++ is written. \textbf{C++14, C++17, and C++23 are \emph{minor} releases} in the formal sense, meaning they introduce no single feature on the scale of move semantics or concepts. But ``minor'' badly understates C++23. It is the release that takes the half-finished cathedral of C++20 and installs the doors, the windows, and the plumbing.

Consider the state of C++20 as actually shipped:

\begin{itemize}
    \item \textbf{Coroutines} were standardized as a \emph{language} mechanism, but the standard library shipped \emph{zero} coroutine types. You could write \lstinline|co_yield|, but you had to hand-roll the promise type, the generator, and the iterator glue yourself. C++23 ships \lstinline|std::generator|.
    \item \textbf{Ranges} shipped with a foundational set of views and algorithms, but the everyday tools --- \lstinline|zip|, \lstinline|enumerate|, \lstinline|chunk|, \lstinline|slide|, fold algorithms, and the all-important \lstinline|ranges::to| for materializing a view into a container --- were cut for time. C++23 ships them.
    \item \textbf{\lstinline|std::format|} gave us a type-safe formatting engine, but no convenient way to actually \emph{write} the result to a stream or the console. You still reached for \lstinline|std::cout << std::format(...)|. C++23 ships \lstinline|std::print| and \lstinline|std::println|.
\end{itemize}

A completion release is one whose value is measured not by what it adds but by what it \emph{finishes}. The features below are deliberately unglamorous: each one removes a workaround you have been carrying since 2020.

\section{The Three Debts C++23 Repays}

It helps to group the work of C++23 around three outstanding debts.

\subsection{The Ergonomics Debt}

Modern C++ had accumulated boilerplate that everyone tolerated because the alternatives were worse. \textbf{Deducing \lstinline|this|} (Chapter~55) collapses the four-overload ref-qualifier explosion and the CRTP boilerplate into a single explicit object parameter. \textbf{\lstinline|std::print|} (Chapter~61) ends the \lstinline|cout|-versus-\lstinline|printf| compromise. \textbf{\lstinline|import std;|} (Chapter~68) replaces dozens of \lstinline|#include| lines with one import. \textbf{\lstinline|size_t| literals} and \textbf{multidimensional \lstinline|operator[]|} (Chapter~66) remove papercuts that have nagged numeric code for a decade.

\subsection{The Error-Handling Debt}

C++ has had two unsatisfying options for reporting failure: exceptions (unpredictable control flow, banned in many low-latency and embedded codebases) and sentinel/\lstinline|optional| returns (which discard \emph{why} the operation failed). \textbf{\lstinline|std::expected<T, E>|} (Chapter~56) gives a value-based channel that carries the error reason, and the \textbf{monadic operations} --- \lstinline|and_then|, \lstinline|transform|, \lstinline|or_else| --- let you compose fallible operations without the nested-\lstinline|if| ``pyramid of doom.'' This is arguably the most consequential library addition in the release for systems programmers.

\subsection{The Standard-Library-Completeness Debt}

The largest body of work is simply \emph{finishing the libraries C++20 started}. \lstinline|std::generator| completes coroutines (Chapter~58). The new range views and fold algorithms complete ranges (Chapters~59--60). Flat containers (Chapter~62), \lstinline|std::stacktrace| (Chapter~63), \lstinline|std::move_only_function| (Chapter~64), and the extended floating-point types (Chapter~67) fill long-standing gaps that forced teams onto Boost, Abseil, or in-house code.

\section{What C++23 Is \emph{Not}}

It is as important to know the boundaries of this release as its contents.

\begin{itemize}
    \item \textbf{C++23 is not reflection.} Static reflection and code injection were targeted but did not land; they are C++26 material and are covered in the next volume.
    \item \textbf{C++23 is not \lstinline|std::execution| (senders/receivers).} The structured-concurrency framework slipped to C++26.
    \item \textbf{C++23 is not contracts.} Contract programming was repeatedly deferred and is, again, C++26.
    \item \textbf{C++23 is not pattern matching.} The \lstinline|inspect| proposal did not make it.
\end{itemize}

Throughout this volume you will see \textbf{version-trap callouts} flagging features that look like they belong to C++23 but are actually C++26 --- a discipline that matters enormously when you target a specific \lstinline|-std| flag in production. The source folder for this volume contained a ``next frontier'' chapter describing exactly these C++26 features; that material is deliberately left out of this volume and reserved for the C++26 volume.

\section{Why This Matters for Low-Latency and Systems Work}

For the high-frequency trading, kernel, and large-scale-distributed audiences this book targets, C++23 is unusually well-aligned with your constraints:

\begin{itemize}
    \item \textbf{\lstinline|std::expected|} gives you exception-free error propagation with zero hidden allocation and a deterministic control-flow cost --- exactly what a hot path or a \lstinline|-fno-exceptions| build needs.
    \item \textbf{Flat containers} trade pointer-chasing trees for cache-friendly contiguous storage, the same transformation you have been doing by hand with sorted \lstinline|vector|s.
    \item \textbf{\lstinline|std::mdspan|} (Chapter~57) gives you a zero-overhead, non-owning, layout-aware multidimensional view --- the missing primitive for numerical kernels that previously demanded a third-party tensor library.
    \item \textbf{\lstinline|std::print|} writes directly to the underlying file descriptor without constructing an intermediate \lstinline|std::string|, making formatted logging cheaper.
    \item \textbf{Expanded \lstinline|constexpr|} (Chapter~65) moves more work to compile time, including \lstinline|constexpr| \lstinline|std::unique_ptr| and a swath of \lstinline|<cmath>|.
\end{itemize}

The recurring theme: C++23 lets you delete the workaround and keep the performance.

\section{How This Volume Is Organized}

The chapters are sequenced from highest-impact to most-specialized:

\begin{itemize}
    \item \textbf{Chapter 55:} Deducing \lstinline|this| --- the marquee core-language feature.
    \item \textbf{Chapters 56--58:} The big library completions: \lstinline|expected|, \lstinline|mdspan|, \lstinline|generator|.
    \item \textbf{Chapters 59--60:} Completing the ranges library: new views, then folds/\lstinline|ranges::to|/search.
    \item \textbf{Chapters 61--62:} Modern output (\lstinline|print|) and flat containers.
    \item \textbf{Chapters 63--64:} Diagnostics, lifetime utilities, and functional/type utilities.
    \item \textbf{Chapters 65--66:} Compile-time refinements and the remaining core-language conveniences.
    \item \textbf{Chapters 67--68:} Extended floating-point types and standard-library modules.
\end{itemize}

Every chapter follows the same structure: motivation, mechanics, edge cases and pitfalls, a low-latency/performance perspective, and complete worked code. Where C++23 merely finishes a C++20 feature, we recap just enough of the C++20 baseline to make the delta clear --- the full treatment lives in Volume 5.

\section{Compiler and Library Support in 2026}

As of this writing, C++23 is broadly usable in production with attention to a few gaps:

\begin{itemize}
    \item \textbf{Core language} features (deducing \lstinline|this|, \lstinline|if consteval|, multidimensional subscript, \lstinline|size_t| literals) are shipping in GCC~13+, Clang~17+, and MSVC~19.3x.
    \item \textbf{\lstinline|std::expected|}, \textbf{\lstinline|std::print|}, \textbf{\lstinline|std::mdspan|}, and the \textbf{new ranges views/algorithms} are available across the three major standard libraries, with \lstinline|std::print|'s direct-to-FD fast path and \lstinline|std::generator| arriving slightly later in libc++.
    \item \textbf{\lstinline|import std;|} has the most uneven support and the strongest build-system dependency; treat it as the one feature to validate carefully on your toolchain before adopting (Chapter~68).
    \item \textbf{Extended floating-point types} (\lstinline|<stdfloat>|) depend on hardware/ABI support and are the most platform-conditional addition (Chapter~67).
\end{itemize}

The practical guidance: compile with \lstinline|-std=c++23| (or \lstinline|/std:c++latest| on MSVC), and gate the two riskiest features --- \lstinline|import std;| and \lstinline|<stdfloat>| --- behind feature-test macros (\lstinline|__cpp_lib_modules|, \lstinline|__STDCPP_FLOAT32_T__|, etc.).

\section{Professional Insights}

\textbf{Treat C++23 as the release that lets you delete code, not add it.} The highest-leverage way to adopt C++23 is to grep your codebase for the workarounds it obsoletes --- hand-rolled generators, the CRTP base classes, the \lstinline|cout << format| idiom, the four-overload accessor pattern, the sorted-\lstinline|vector|-as-map --- and replace them with the standard facility. Each replacement removes a maintenance liability and usually improves performance.

\textbf{Adopt \lstinline|std::expected| first; it changes how you design APIs.} Of everything in this release, \lstinline|expected| has the deepest ripple effect, because it lets you make ``this operation can fail'' a visible part of a function's type rather than a runtime surprise. Codebases that ban exceptions gain a principled error channel for the first time; codebases that use exceptions gain a cheaper one for expected failures. Plan for it to propagate through your interfaces.

\textbf{Pin the standard and audit feature-test macros, because ``C++23'' is not monolithic.} Different features matured on different timelines across the three toolchains. Do not assume that because deducing \lstinline|this| compiles, \lstinline|import std;| and \lstinline|<stdfloat>| will too. Use \lstinline|__has_include|, the \lstinline|__cpp_*| feature-test macros, and CI matrix builds across your target compilers to verify exactly which subset of C++23 you can rely on.

\textbf{Keep the C++26 line bright.} Several of the most-hyped ``modern C++'' features people associate with this era --- reflection, senders/receivers, contracts, pattern matching --- are \emph{not} in C++23. Knowing precisely where the C++23 boundary lies prevents you from writing code that fails to compile under the \lstinline|-std=c++23| flag your build system actually passes.
